\section{PPXA --- proximal splitting twin of Adam}
\label{sec:ppxa}

\paragraph{How it works.}
PPXA minimises the \emph{same} objective~\eqref{eq:objective} as Adam, but by evaluating each
term's proximal operator in parallel, averaging, and over-relaxing:
\begin{quote}
$p_1 \leftarrow (I+\gamma H^{\!\top}\!H)^{-1}(u_1+\gamma H^{\!\top} g)$ \hfill\textit{(data prox: pull toward $g$)}\\
$p_2 \leftarrow \mathrm{prox}_{\gamma\lambda R}(u_2)$ \hfill\textit{(prior prox)}\\
$p_3 \leftarrow \max(u_3,0)$ \hfill\textit{(positivity prox)}\\
$\bar p \leftarrow \tfrac13(p_1+p_2+p_3)$ \hfill\textit{(average)}\\
$u_i \leftarrow u_i + \rho_{\mathrm{rel}}(2\bar p - f - p_i)$;\quad $f \leftarrow f + \rho_{\mathrm{rel}}(\bar p - f)$ \hfill\textit{(over-relax)}
\end{quote}
\noindent where $\gamma$ is the prox step (gain $\gamma s_i^2$ per singular direction), $u_1,u_2,u_3$
the auxiliary variables, $R$ the regulariser at share $\lambda$, and $\rho_{\mathrm{rel}}\in(0,2)$
the relaxation. Minimising the same convex $L$, it reaches the \emph{same image} as Adam.

\paragraph{Key parameters.}
The same solution axes as Adam (\texttt{reg}, $\lambda$); its own knobs set speed only, provided the
proxes are exact: $\gamma\in[1/s_1^2,1/s_2^2]$ (out of range it can stall) and
$\rho_{\mathrm{rel}}=1.5$.

\paragraph{Results.}
Run identically to Adam. The campaign confirms the framework claim --- \emph{Adam and PPXA return
the same image} on the shared L1 objective (reference structure, noise 0.02):
\IfFileExists{tables/algo_ppxa.tex}{\paragraph{matirf}
\begin{center}\begin{adjustbox}{max width=\linewidth}
\begin{tabular}{lllllllll}
\toprule
use & dataset & noise & params & NMSE & PSNR & Depth Error (nm) & Stack Recovery & chi2\_ratio \\
\midrule
optimal & esoubies & native & reg=l1, lambda=0.2 & — & — & — & — & 18.8 \\
optimal & fibres & 0.02 & reg=l1, lambda=0.2 & 0.855 & 24.7 & 73 & 0 & 0.446 \\
\bottomrule
\end{tabular}\end{adjustbox}\end{center}
}{}

\paragraph{Limit.}
PPXA shares Adam's ceiling (only as good as the analytic prior) and adds two costs: it needs
\emph{proximable} priors (Adam takes any differentiable one), and its splitting is markedly slower
here (median runtime several times Adam's).
